\scanpage{368}
\mspage{92}
\unsure{6.5.|P.S.} Finalement, la relation entrevue avec les ``types d'homotopie simples'' de Whitehead s'évanouit\textsuperscript{*} : en effet, ces derniers sont strictement
% slanted
\margin{Mais peut-être quand même, -- la question reste ouverte.} plus fins que les types d'homotopie (\uline{stricts}) ordinaires, et l'invariant de Whitehead n'est \uline{pas} un invariant du type d'homotopie ordinaire. Il est vrai que lorsqu'on considère le squelette simplicial d'un polyèdre triangulé comme objet de $\mathrm{Cat}$, l'homotopie (continue) ordinaire entre deux tels polyèdres s'interprète, dans $\mathrm{Cat}$, non pas \deleted{\unsure{(contr.}} (me semble-t-il) par l'existence d'un homotopisme\deleted{\unsure{,}} entre les objets correspondants de $\mathrm{Cat}$, mais par leur isomorphie dans $\mathrm{Hot}_{\mathbb W_\infty}\overset{\mathrm{def}}{=}\mathrm{Cat}(\mathbb W_\infty)^{-1}$. Mais le procédé d'approximations simpliciales de Brouwer montre que si les polyèdres $X$ et $Y$ sont homotopes, alors il existe des subdivisions simpliciales \inserted{$X',Y'$} de $X$ et de $Y$ qui sont ``combinatoirement homotopes'' (dans un sens convenable). Donc \deleted{question qui \unsure{s'apprête|s'impose}} je présume que les maquettes simpliciales de $X'$ et de $Y'$ sont alors homotopes dans $\mathrm{Cat}$. Il s'ensuivrait que
% archive sheet number: 367
